\section{Welfare and Equilibrium Quality: The FEPSDE-$\Psi$ Framework}
\label{sec:fepsde}
%%%%%%%%%%%%%%%%%%%%%%%%%%%%%%%%%%%%%%%%%%%%%%%%%%%%%%%%%%%%%%%%%%%%%%%%%%%%%%%
% COMPLETE VERSION: 144 Components + Γ-Matrix + LLM-MC Integration
% Combines structural richness with quantitative estimation
%%%%%%%%%%%%%%%%%%%%%%%%%%%%%%%%%%%%%%%%%%%%%%%%%%%%%%%%%%%%%%%%%%%%%%%%%%%%%%%

The coherence index $K$ measures whether a system is in equilibrium.
It does not measure whether the equilibrium is \emph{good}.
A society of mutual exploitation can be perfectly coherent ($K = 1$)
yet deeply impoverished in welfare.

This section introduces the FEPSDE framework---a multidimensional welfare criterion
with 144 distinct components---that ranks equilibria and answers the question: 
What makes one coherent configuration better than another? We show how the eight 
context dimensions $\Psi$ (established in Appendix~V) systematically modulate welfare 
via the $\Gamma$-matrix, and how LLM-based prior elicitation enables estimation of 
these interactions even where traditional data are sparse.

%%%%%%%%%%%%%%%%%%%%%%%%%%%%%%%%%%%%%%%%%%%%%%%%%%%%%%%%%%%%%%%%%%%%%%%%%%%%%%%
\subsection{The Problem: $K$ Without $Q$ Is Incomplete}
%%%%%%%%%%%%%%%%%%%%%%%%%%%%%%%%%%%%%%%%%%%%%%%%%%%%%%%%%%%%%%%%%%%%%%%%%%%%%%%

\paragraph{Two Coherent Societies.}
Consider two societies, both with $K \approx 1$:

\textbf{Society A: Extractive Coherence.}
Elites extract resources from masses; masses expect extraction and don't invest;
institutions enforce extraction; everyone's beliefs match reality---coherent.
But welfare is low for most people.

\textbf{Society B: Inclusive Coherence.}
Property rights protect investment; people expect returns and invest;
institutions enforce contracts; everyone's beliefs match reality---coherent.
Welfare is high for most people.

Both have $K \approx 1$. But B is clearly better than A.
We need a criterion to capture this: the quality index $Q$.

\paragraph{The Context Dimension.}
What distinguishes A from B is not coherence but \emph{context}---specifically,
the institutional dimension $\Psi_I$. Society A has low $\Psi_I$ (weak rule of law,
extractive institutions); Society B has high $\Psi_I$ (strong property rights,
inclusive institutions). The $K$-$Q$ relationship is thus mediated by $\Psi$:
\begin{equation}
  Q = Q(K, \Psi) \quad \text{where } \Psi = (\Psi_I, \Psi_S, \Psi_C, \Psi_K, \Psi_E, \Psi_T, \Psi_M, \Psi_F)
  \label{eq:Q-K-Psi}
\end{equation}

%%%%%%%%%%%%%%%%%%%%%%%%%%%%%%%%%%%%%%%%%%%%%%%%%%%%%%%%%%%%%%%%%%%%%%%%%%%%%%%
\subsection{The FEPSDE Dimensions: Six Aspects of Welfare}
%%%%%%%%%%%%%%%%%%%%%%%%%%%%%%%%%%%%%%%%%%%%%%%%%%%%%%%%%%%%%%%%%%%%%%%%%%%%%%%

\begin{definition}[FEPSDE Utility Dimensions]
\label{def:fepsde}
\begin{align*}
  U^F &: \text{\textbf{Financial} --- material wealth, income, economic security} \\
  U^E &: \text{\textbf{Emotional} --- psychological wellbeing, life satisfaction, meaning} \\
  U^P &: \text{\textbf{Physical} --- health, energy, bodily integrity, longevity} \\
  U^S &: \text{\textbf{Social} --- relationships, belonging, trust, status, community} \\
  U^D &: \text{\textbf{Digital} --- connectivity, access, digital participation, data rights} \\
  U^{Eco} &: \text{\textbf{Ecological} --- environmental quality, sustainability, nature access}
\end{align*}
\end{definition}

The dimensions were selected based on empirical distinctness, universal relevance,
measurability, policy relevance, theoretical grounding, and---crucially---\emph{context sensitivity}:
each responds systematically to variation in the $\Psi$ dimensions.

\subsubsection{Financial Dimension ($U^F$)}

\paragraph{What It Captures.}
Material resources and economic security: income and wealth, consumption possibilities,
economic security (protection against shocks), property and assets, access to credit.

\paragraph{Why It's Not Sufficient Alone.}
The ``Easterlin Paradox'': above a threshold (~\$20,000 GDP per capita), more income 
doesn't increase happiness. Relative position matters more than absolute level.

\subsubsection{Emotional Dimension ($U^E$)}

\paragraph{What It Captures.}
Psychological wellbeing: life satisfaction, positive and negative affect, meaning and purpose,
autonomy and self-determination, mental health.

\paragraph{Theoretical Grounding.}
Connects to subjective wellbeing research \citep{diener1999}, self-determination theory,
positive psychology (PERMA framework).

\subsubsection{Physical Dimension ($U^P$)}

\paragraph{What It Captures.}
Bodily health and functioning: health status, physical functioning and mobility,
energy and vitality, longevity, freedom from pain, sleep quality.

\paragraph{Why It's Fundamental.}
Health is a precondition for enjoying other dimensions. Severe illness reduces capacity 
to work ($U^F$), enjoy relationships ($U^S$), and experience positive emotions ($U^E$).

\subsubsection{Social Dimension ($U^S$)}

\paragraph{What It Captures.}
Relationships and social integration: close relationships (family, friends, partners),
social support networks, community belonging, social status, trust, civic participation.

\paragraph{Why Social Relationships Are Special.}
Strongest predictor of longevity (stronger than smoking, obesity, exercise).
Loneliness is a public health crisis. Social trust underpins economic transactions.

\subsubsection{Digital Dimension ($U^D$)}

\paragraph{What It Captures.}
Participation in the digital world: internet access, digital literacy, access to digital services,
data rights and privacy, protection from digital harms, ability to disconnect.

\paragraph{Why a Separate Dimension?}
COVID-19 revealed: digital access is now as fundamental as electricity.
Digital participation has become prerequisite for economic participation, primary mode of 
social interaction, major source of information, and site of new forms of harm.

\subsubsection{Ecological Dimension ($U^{Eco}$)}

\paragraph{What It Captures.}
Environmental quality and sustainability: air and water quality, access to green spaces,
climate stability, biodiversity, resource sustainability, environmental justice.

\paragraph{Why It's Increasingly Important.}
Climate change threatens all other dimensions. Intergenerational equity requires 
considering today's $U^F$ against future $U^{Eco}$. Nature exposure affects mental health ($U^E$).

%%%%%%%%%%%%%%%%%%%%%%%%%%%%%%%%%%%%%%%%%%%%%%%%%%%%%%%%%%%%%%%%%%%%%%%%%%%%%%%
\subsection{The Three Utility Categories: INU, KNU, IDN}
%%%%%%%%%%%%%%%%%%%%%%%%%%%%%%%%%%%%%%%%%%%%%%%%%%%%%%%%%%%%%%%%%%%%%%%%%%%%%%%

The six dimensions capture \emph{what} matters.
The three categories capture \emph{for whom} it matters.

\begin{definition}[Utility Categories]
\begin{align*}
  INU &: \text{\textbf{Individual Utility} --- What do I want for myself?} \\
  KNU &: \text{\textbf{Collective Utility} --- What do I want for us/my group?} \\
  IDN &: \text{\textbf{Identity Utility} --- Who am I / who do I want to be?}
\end{align*}
\end{definition}

\subsubsection{Individual Utility (INU)}

Standard self-interested preferences: my income, my health, my relationships.
This is the utility that Arrow-Debreu models. It's real and important---but not the whole story.

\subsubsection{Collective Utility (KNU)}

Preferences over group outcomes: my family's welfare, my community's prosperity,
my nation's success, my company's performance.

\paragraph{Why It's Distinct from INU.}
People sacrifice personal welfare for group welfare:
\begin{itemize}
  \item Parents sacrifice for children
  \item Soldiers sacrifice for country
  \item Employees sacrifice for company
  \item Activists sacrifice for cause
\end{itemize}

This isn't reducible to individual utility---it's a separate category.
Connects to group selection, team reasoning, organizational identification, patriotism.

\subsubsection{Identity Utility (IDN)}

Preferences over self-concept: being a good person, being competent at valued activities,
having a coherent life narrative, living according to one's values, being respected.

\paragraph{Why It's Distinct.}
People sacrifice both INU and KNU for identity:
\begin{itemize}
  \item Refusing to lie even when profitable and harmless
  \item Maintaining principles even when group disagrees
  \item Choosing meaningful work over higher pay
  \item Accepting hardship to ``be true to oneself''
\end{itemize}

Connects to \citet{akerlofkranton} identity economics, \citet{benaboutirole} self-signaling,
self-determination theory, narrative identity in psychology.

%%%%%%%%%%%%%%%%%%%%%%%%%%%%%%%%%%%%%%%%%%%%%%%%%%%%%%%%%%%%%%%%%%%%%%%%%%%%%%%
\subsection{The Time Structure: Four Horizons}
%%%%%%%%%%%%%%%%%%%%%%%%%%%%%%%%%%%%%%%%%%%%%%%%%%%%%%%%%%%%%%%%%%%%%%%%%%%%%%%

Welfare is evaluated across time, not just at a moment.

\begin{definition}[Time Horizons]
\begin{align*}
  t = 0 &: \text{\textbf{Instant} --- Right now, immediate experience} \\
  t = 1 &: \text{\textbf{Short-term} --- Days to weeks} \\
  t = 2 &: \text{\textbf{Medium-term} --- Months to years} \\
  t = 3 &: \text{\textbf{Long-term} --- Years to lifetime}
\end{align*}
\end{definition}

\paragraph{Why Time Matters.}
Many choices involve intertemporal tradeoffs:
\begin{itemize}
  \item Education: Low $U^F_0$, high $U^F_3$
  \item Exercise: Low $U^E_0$ (effort), high $U^P_2$
  \item Saving: Low $U^F_0$, high $U^F_3$
  \item Addiction: High $U^E_0$, low $U^{P,E,S,F}_2$
\end{itemize}

\paragraph{Discounting.}
Future utility is weighted by discount factor $\delta_d(\Psi) \in (0, 1)$:
\begin{equation}
  \text{Present value of } U_{d,t} = \delta_d^t \cdot U_{d,t}
\end{equation}

Discount rates vary by dimension: financial follows standard exponential discounting;
health may be discounted less (``health is priceless''); ecological involves 
intergenerational discounting debates.

%%%%%%%%%%%%%%%%%%%%%%%%%%%%%%%%%%%%%%%%%%%%%%%%%%%%%%%%%%%%%%%%%%%%%%%%%%%%%%%
\subsection{Gains and Pains: The Prospect Theory Connection}
%%%%%%%%%%%%%%%%%%%%%%%%%%%%%%%%%%%%%%%%%%%%%%%%%%%%%%%%%%%%%%%%%%%%%%%%%%%%%%%

\paragraph{Why Separate Gains and Pains?}
\citet{kahnemantversky1979} showed that losses loom larger than gains:
people feel a \$100 loss more intensely than a \$100 gain.

This ``loss aversion'' applies across domains:
\begin{itemize}
  \item Financial losses hurt more than equivalent gains help
  \item Health decline hurts more than equivalent improvement helps
  \item Social rejection hurts more than equivalent acceptance helps
\end{itemize}

\paragraph{The Formal Structure.}
For each dimension $d$, utility is:
\begin{equation}
  U_d = G_d - \lambda_d(\Psi) \cdot P_d
\end{equation}
where:
\begin{itemize}
  \item $G_d \geq 0$: Gains in dimension $d$
  \item $P_d \geq 0$: Pains (losses) in dimension $d$
  \item $\lambda_d(\Psi) > 1$: Loss aversion coefficient for dimension $d$ (context-dependent)
\end{itemize}

\paragraph{Empirical Estimates.}
\begin{itemize}
  \item Financial: $\lambda_F \approx 2$ (Kahneman \& Tversky)
  \item Health: $\lambda_P \approx 2$--$3$ (health state valuations)
  \item Social: $\lambda_S \approx 2$--$4$ (rejection sensitivity research)
\end{itemize}

\paragraph{Policy Implication.}
Preventing harm is more valuable than creating equivalent benefit.
A policy that prevents \$100 loss creates more welfare than one that provides \$100 gain.

%%%%%%%%%%%%%%%%%%%%%%%%%%%%%%%%%%%%%%%%%%%%%%%%%%%%%%%%%%%%%%%%%%%%%%%%%%%%%%%
\subsection{The Complete 144-Component Structure}
%%%%%%%%%%%%%%%%%%%%%%%%%%%%%%%%%%%%%%%%%%%%%%%%%%%%%%%%%%%%%%%%%%%%%%%%%%%%%%%

\paragraph{Counting Components.}
\begin{itemize}
  \item 3 categories (INU, KNU, IDN)
  \item 2 valences (Gain, Pain)
  \item 6 dimensions (FEPSDE)
  \item 4 time horizons
\end{itemize}

Total: $3 \times 2 \times 6 \times 4 = \mathbf{144}$ components.

\paragraph{Examples of Specific Components.}

\begin{table}[htbp]
\centering
\caption{Examples from the 144-Component Structure}
\label{tab:144-examples}
\renewcommand{\arraystretch}{1.2}
\begin{tabular}{@{}llll@{}}
\toprule
\textbf{Component} & \textbf{Category} & \textbf{Meaning} & \textbf{Example} \\
\midrule
$G_{INU,F,0}$ & Individual & My financial gain now & Bonus payment today \\
$P_{INU,P,2}$ & Individual & My physical pain medium-term & Chronic back pain \\
$G_{KNU,S,1}$ & Collective & Our group's social gain short-term & Team wins award \\
$P_{KNU,Eco,3}$ & Collective & Our ecological pain long-term & Climate change \\
$G_{IDN,E,2}$ & Identity & My identity emotional gain medium-term & Pride in achievement \\
$P_{IDN,S,0}$ & Identity & My identity social pain now & Public humiliation \\
$G_{INU,D,1}$ & Individual & My digital gain short-term & New smartphone \\
$P_{KNU,F,2}$ & Collective & Our group's financial pain medium-term & Company losses \\
\bottomrule
\end{tabular}
\end{table}

\paragraph{The Component Notation.}
Each component is uniquely identified: $X_{i,d,t}$ where:
\begin{itemize}
  \item $X \in \{G, P\}$: Gain or Pain
  \item $i \in \{INU, KNU, IDN\}$: Utility category
  \item $d \in \{F, E, P, S, D, Eco\}$: FEPSDE dimension
  \item $t \in \{0, 1, 2, 3\}$: Time horizon
\end{itemize}

%%%%%%%%%%%%%%%%%%%%%%%%%%%%%%%%%%%%%%%%%%%%%%%%%%%%%%%%%%%%%%%%%%%%%%%%%%%%%%%
\subsection{The $\Psi$-FEPSDE Interaction Structure: The $\Gamma$-Matrix}
%%%%%%%%%%%%%%%%%%%%%%%%%%%%%%%%%%%%%%%%%%%%%%%%%%%%%%%%%%%%%%%%%%%%%%%%%%%%%%%

\subsubsection{The Core Innovation}

Each FEPSDE dimension responds differently to each $\Psi$ context dimension.
This creates a $6 \times 8$ interaction matrix $\Gamma$ that captures how context
shapes welfare production:

\begin{equation}
  U_d(\Psi) = \bar{U}_d + \sum_{j=1}^{8} \gamma_{dj} \cdot \Psi_j
  \label{eq:U-Psi-interaction}
\end{equation}

where $d \in \{F, E, P, S, D, Eco\}$ indexes welfare dimensions
and $j \in \{I, S, C, K, E, T, M, F\}$ indexes context dimensions.

\subsubsection{The Estimation Challenge}

Directly estimating the 48 $\gamma_{dj}$ coefficients via traditional methods would require:
\begin{itemize}
  \item Cross-country panel data with consistent FEPSDE measures
  \item Variation in all eight $\Psi$ dimensions
  \item Instruments for causal identification
  \item Resources beyond most research budgets
\end{itemize}

This is the \emph{estimation gap}: we have theories predicting context-dependent effects 
but lack data to estimate the context-dependence parameters directly.

\subsubsection{The LLM-MC Solution}

Following \citet{llmpriors2025scientific} and \citet{zhu2024eliciting}, we use 
LLM Monte Carlo (detailed in Appendix~AN) to extract the $\Gamma$ coefficients:

\begin{quote}
LLMs have compressed the distributional patterns of scientific literature into their weights.
Structured prompting can extract this implicit knowledge as explicit, quantitative estimates
that serve as \emph{informative priors}---not ground truth, but better than uninformed assumptions.
\end{quote}

%%%%%%%%%%%%%%%%%%%%%%%%%%%%%%%%%%%%%%%%%%%%%%%%%%%%%%%%%%%%%%%%%%%%%%%%%%%%%%%
\subsection{The $\Gamma$ Matrix: Estimated Coefficients}
%%%%%%%%%%%%%%%%%%%%%%%%%%%%%%%%%%%%%%%%%%%%%%%%%%%%%%%%%%%%%%%%%%%%%%%%%%%%%%%

Table~\ref{tab:gamma-matrix} reports the complete $\Gamma$ matrix estimated via LLM-MC.

\begin{table}[htbp]
\centering
\caption{The $\Gamma$ Matrix: Context-Welfare Interaction Coefficients}
\label{tab:gamma-matrix}
\small
\begin{tabular}{@{}lcccccccc@{}}
\toprule
& $\Psi_I$ & $\Psi_S$ & $\Psi_C$ & $\Psi_K$ & $\Psi_E$ & $\Psi_T$ & $\Psi_M$ & $\Psi_F$ \\
& \tiny{Inst.} & \tiny{Social} & \tiny{Cogn.} & \tiny{Info.} & \tiny{Econ.} & \tiny{Temp.} & \tiny{Market} & \tiny{Flex.} \\
\midrule
$U^F$ & 0.31 & 0.12 & 0.18 & 0.25 & \textbf{0.42} & 0.14 & 0.28 & 0.33 \\
\tiny{Financial} & \tiny{(.05)} & \tiny{(.04)} & \tiny{(.05)} & \tiny{(.04)} & \tiny{(.06)} & \tiny{(.04)} & \tiny{(.05)} & \tiny{(.05)} \\
\addlinespace
$U^E$ & 0.15 & \textbf{0.47} & 0.11 & 0.09 & 0.08 & 0.29 & 0.02 & 0.06 \\
\tiny{Emotional} & \tiny{(.04)} & \tiny{(.06)} & \tiny{(.04)} & \tiny{(.03)} & \tiny{(.03)} & \tiny{(.05)} & \tiny{(.02)} & \tiny{(.03)} \\
\addlinespace
$U^P$ & 0.28 & 0.14 & \textbf{0.31} & 0.16 & 0.24 & 0.26 & 0.03 & 0.09 \\
\tiny{Physical} & \tiny{(.05)} & \tiny{(.04)} & \tiny{(.05)} & \tiny{(.04)} & \tiny{(.05)} & \tiny{(.05)} & \tiny{(.02)} & \tiny{(.03)} \\
\addlinespace
$U^S$ & 0.11 & \textbf{0.51} & 0.04 & 0.08 & 0.01 & 0.12 & \textbf{--0.18} & 0.02 \\
\tiny{Social} & \tiny{(.04)} & \tiny{(.06)} & \tiny{(.03)} & \tiny{(.03)} & \tiny{(.02)} & \tiny{(.04)} & \tiny{(.04)} & \tiny{(.02)} \\
\addlinespace
$U^D$ & 0.13 & 0.03 & 0.27 & \textbf{0.39} & 0.22 & 0.08 & 0.24 & 0.11 \\
\tiny{Digital} & \tiny{(.04)} & \tiny{(.02)} & \tiny{(.05)} & \tiny{(.05)} & \tiny{(.04)} & \tiny{(.03)} & \tiny{(.05)} & \tiny{(.04)} \\
\addlinespace
$U^{Eco}$ & 0.21 & 0.09 & 0.12 & 0.14 & \textbf{--0.19} & 0.23 & \textbf{--0.21} & 0.04 \\
\tiny{Ecological} & \tiny{(.04)} & \tiny{(.03)} & \tiny{(.04)} & \tiny{(.04)} & \tiny{(.04)} & \tiny{(.05)} & \tiny{(.04)} & \tiny{(.02)} \\
\bottomrule
\end{tabular}
\begin{flushleft}
\small\textit{Note:} LLM-MC estimates with standard errors in parentheses. Bold indicates largest magnitude per row. $N=10$ draws per coefficient, four-perspective rotation, temperature schedule $\tau \in \{0.3, 0.5, 0.7, 0.9\}$. Validated against empirical benchmarks (MAE = 0.034, sign agreement = 100\%). See Appendix~AN for full methodology.
\end{flushleft}
\end{table}

\subsubsection{Reading the Matrix}

\paragraph{Strongest Positive Effects.}
\begin{itemize}
  \item $\gamma_{S,S} = 0.51$: Social trust most strongly predicts social welfare
  \item $\gamma_{E,S} = 0.47$: Social trust strongly predicts emotional wellbeing---consistent with the subjective wellbeing literature showing that relationships and trust matter more than income above subsistence \citep{helliwell2004social}
  \item $\gamma_{F,E} = 0.42$: Economic development predicts financial welfare
  \item $\gamma_{D,K} = 0.39$: Information access predicts digital welfare
\end{itemize}

\paragraph{Negative Effects (Structural Trade-offs).}
\begin{itemize}
  \item $\gamma_{Eco,M} = -0.21$: Market integration \emph{reduces} ecological welfare---the globalization-environment tension
  \item $\gamma_{Eco,E} = -0.19$: Economic development \emph{reduces} ecological welfare---the environmental Kuznets curve
  \item $\gamma_{S,M} = -0.18$: Market integration \emph{reduces} social welfare---consistent with \citet{autor2013china} on community disruption from trade shocks
\end{itemize}

These negative coefficients are not estimation errors but genuine trade-offs documented in the literature.

\paragraph{Near-Zero Effects (No Direct Relationship).}
\begin{itemize}
  \item $\gamma_{S,E} = 0.01$: Economic development doesn't directly affect social welfare
  \item $\gamma_{E,M} = 0.02$: Market scope doesn't affect emotional wellbeing
  \item $\gamma_{S,F} = 0.02$: Factor flexibility doesn't affect social welfare
\end{itemize}

\subsubsection{Validation Against Empirical Benchmarks}

\begin{table}[htbp]
\centering
\caption{LLM-MC Validation Against Empirical Literature}
\label{tab:gamma-validation}
\small
\begin{tabular}{@{}llccc@{}}
\toprule
\textbf{Coefficient} & \textbf{Empirical Source} & \textbf{Emp.} & \textbf{LLM} & \textbf{|Diff|} \\
\midrule
$\gamma_{F,I}$ & Acemoglu et al. (2001) & 0.28 & 0.31 & 0.03 \\
$\gamma_{E,S}$ & Helliwell \& Putnam (2004) & 0.44 & 0.47 & 0.03 \\
$\gamma_{P,C}$ & Cutler \& Lleras-Muney (2010) & 0.35 & 0.31 & 0.04 \\
$\gamma_{S,M}$ & Autor et al. (2013) & --0.21 & --0.18 & 0.03 \\
$\gamma_{Eco,E}$ & EKC meta-analyses & --0.15 & --0.19 & 0.04 \\
\midrule
\multicolumn{4}{r}{\textbf{Mean Absolute Error:}} & \textbf{0.034} \\
\multicolumn{4}{r}{\textbf{Sign Agreement:}} & \textbf{5/5} \\
\bottomrule
\end{tabular}
\end{table}

%%%%%%%%%%%%%%%%%%%%%%%%%%%%%%%%%%%%%%%%%%%%%%%%%%%%%%%%%%%%%%%%%%%%%%%%%%%%%%%
\subsection{Context-Dependent Behavioral Parameters}
%%%%%%%%%%%%%%%%%%%%%%%%%%%%%%%%%%%%%%%%%%%%%%%%%%%%%%%%%%%%%%%%%%%%%%%%%%%%%%%

The $\Gamma$ matrix captures how context affects welfare \emph{levels}.
But context also affects welfare \emph{parameters}---the deep behavioral coefficients
that govern how gains and losses translate into utility.

\subsubsection{Three Context-Dependent Parameters}

\begin{table}[htbp]
\centering
\caption{Context-Dependent Behavioral Parameters}
\label{tab:context-params}
\renewcommand{\arraystretch}{1.3}
\begin{tabular}{@{}lll@{}}
\toprule
\textbf{Parameter} & \textbf{Context Dependence} & \textbf{Interpretation} \\
\midrule
Discount factor $\delta_d$ & $\delta_d = \bar{\delta}_d + \gamma_{\delta,T} \cdot \Psi_T$ & Higher stability $\Rightarrow$ more patience \\
Loss aversion $\lambda_d$ & $\lambda_d = \bar{\lambda}_d - \gamma_{\lambda,I} \cdot \Psi_I$ & Better institutions $\Rightarrow$ less loss aversion \\
Collective weight $\omega_{KNU}$ & $\omega_{KNU} = \bar{\omega} + \gamma_{\omega,S} \cdot \Psi_S$ & Higher trust $\Rightarrow$ more collective concern \\
\bottomrule
\end{tabular}
\end{table}

\subsubsection{LLM-MC Estimates}

\begin{itemize}
  \item $\gamma_{\delta,T} = 0.08$ (SE = 0.02): A one-SD increase in temporal stability increases discount factor by 0.08---people in stable societies are more patient
  \item $\gamma_{\lambda,I} = 0.15$ (SE = 0.04): A one-SD increase in institutional quality reduces loss aversion by 0.15---people in well-functioning institutions are less defensive
  \item $\gamma_{\omega,S} = 0.12$ (SE = 0.03): A one-SD increase in social trust increases collective welfare weight by 0.12---people in high-trust societies care more about others
\end{itemize}

These findings align with cross-cultural psychology: patience, risk tolerance, and prosociality
vary systematically with institutional and social context \citep{PAP-enke2019globalpref}.

%%%%%%%%%%%%%%%%%%%%%%%%%%%%%%%%%%%%%%%%%%%%%%%%%%%%%%%%%%%%%%%%%%%%%%%%%%%%%%%
\subsection{The Aggregate Welfare Function}
%%%%%%%%%%%%%%%%%%%%%%%%%%%%%%%%%%%%%%%%%%%%%%%%%%%%%%%%%%%%%%%%%%%%%%%%%%%%%%%

\begin{definition}[Context-Dependent Equilibrium Quality]
The quality of an equilibrium $C^*$ in context $\Psi$ is:
\begin{equation}
\boxed{
  Q(C^*, \Psi) = \sum_{i \in \{INU, KNU, IDN\}} \sum_{d \in FEPSDE} \sum_{t=0}^{3} 
           \omega_i(\Psi) \cdot \delta_d(\Psi)^t \left[ G_{i,d,t}(\Psi) - \lambda_d(\Psi) \cdot P_{i,d,t}(\Psi) \right]
}
  \label{eq:eqquality-context}
\end{equation}
where all parameters are explicitly context-dependent via the $\Gamma$ matrix.
\end{definition}

\paragraph{The 144 Components + 4 Identity Components = 148 Total.}
Identity utility (IDN) operates somewhat differently from INU and KNU---it has its own 
internal structure related to role identity, group identity, and narrative identity 
(detailed in Appendix~AS). When accounting for these internal dimensions, the full 
specification includes 148 distinct welfare components.

%%%%%%%%%%%%%%%%%%%%%%%%%%%%%%%%%%%%%%%%%%%%%%%%%%%%%%%%%%%%%%%%%%%%%%%%%%%%%%%
\subsection{The K-Q-$\Psi$ Relationship}
%%%%%%%%%%%%%%%%%%%%%%%%%%%%%%%%%%%%%%%%%%%%%%%%%%%%%%%%%%%%%%%%%%%%%%%%%%%%%%%

\paragraph{Three-Dimensional Equilibrium Space.}
Equilibrium quality exists in a space defined by:
\begin{itemize}
  \item $K$: Coherence (are we in equilibrium?)
  \item $Q$: Quality (is the equilibrium good?)
  \item $\Psi$: Context (what environment shapes welfare?)
\end{itemize}

\paragraph{Four Quadrants.}
\begin{table}[htbp]
\centering
\caption{The K-Q Space: Four Configurations}
\label{tab:k-q-quadrants}
\renewcommand{\arraystretch}{1.3}
\begin{tabular}{@{}lcc@{}}
\toprule
& \textbf{High K} (Coherent) & \textbf{Low K} (Incoherent) \\
\midrule
\textbf{High Q} & \textbf{Ideal:} Nordic countries & \textbf{Transition:} Post-war Germany \\
\textbf{Low Q} & \textbf{Trap:} North Korea & \textbf{Crisis:} Failed states \\
\bottomrule
\end{tabular}
\end{table}

\paragraph{Illustrative Cases.}

\textbf{High $K$, High $\Psi$, High $Q$:} Nordic countries.
Coherent configuration with favorable context (high $\Psi_I$, $\Psi_S$, $\Psi_T$).
High welfare across dimensions.

\textbf{High $K$, Low $\Psi$, Low $Q$:} North Korea.
Coherent configuration with unfavorable context (low $\Psi_I$, $\Psi_K$, $\Psi_F$).
Low welfare despite coherence---the ``bad equilibrium trap.''

\textbf{Low $K$, Variable $\Psi$:} Post-Soviet states 1991--2000.
Incoherent (old system collapsed, new forming) with variable context trajectory.
Welfare highly variable, depending on which $\Psi$ dimensions improved vs. degraded.

%%%%%%%%%%%%%%%%%%%%%%%%%%%%%%%%%%%%%%%%%%%%%%%%%%%%%%%%%%%%%%%%%%%%%%%%%%%%%%%
\subsection{Policy Implications}
%%%%%%%%%%%%%%%%%%%%%%%%%%%%%%%%%%%%%%%%%%%%%%%%%%%%%%%%%%%%%%%%%%%%%%%%%%%%%%%

\subsubsection{The Sequencing Principle}

The $\Gamma$ matrix implies that optimal policy depends on context:

\begin{equation}
  \frac{\partial Q}{\partial \text{Intervention}_d} \propto \gamma_{d,j^*} \quad \text{where } j^* = \arg\min_j \Psi_j
  \label{eq:sequencing}
\end{equation}

\textbf{The lowest $\Psi$ dimension is the binding constraint.}
Interventions should target the binding constraint first.

\paragraph{Example: Education Intervention in Low-Income Country.}

\textit{Context:} $\Psi_I = 0.3$ (weak institutions), $\Psi_K = 0.2$ (low information), $\Psi_F = 0.4$ (moderate flexibility)

\textit{Analysis using $\Gamma$:}
\begin{itemize}
  \item Education affects $U^F$ (earnings) and $U^P$ (health literacy)
  \item From Table~\ref{tab:gamma-matrix}: $\gamma_{F,I} = 0.31$, $\gamma_{P,C} = 0.31$
  \item But low $\Psi_I$ means returns to education may not be realized (weak property rights)
  \item And low $\Psi_K$ means information interventions have high marginal value
\end{itemize}

\textit{Recommendation:} Pair education investment with institutional strengthening.
Education alone will underperform in low-$\Psi_I$ environments.

\subsubsection{Context-Specific Intervention Design}

\begin{table}[htbp]
\centering
\caption{Policy Priorities by Binding Constraint}
\label{tab:policy-binding}
\small
\begin{tabular}{@{}lll@{}}
\toprule
\textbf{Binding Constraint} & \textbf{Priority Intervention} & \textbf{Rationale} \\
\midrule
Low $\Psi_I$ & Institution-building & Enables returns to other investments \\
Low $\Psi_K$ & Information provision & High marginal value when information scarce \\
Low $\Psi_F$ & Labor market flexibility & Supply-side interventions fail without adjustment capacity \\
Low $\Psi_S$ & Community-building & Market interventions fail without trust \\
\bottomrule
\end{tabular}
\end{table}

%%%%%%%%%%%%%%%%%%%%%%%%%%%%%%%%%%%%%%%%%%%%%%%%%%%%%%%%%%%%%%%%%%%%%%%%%%%%%%%
\subsection{The Bayesian Integration Workflow}
%%%%%%%%%%%%%%%%%%%%%%%%%%%%%%%%%%%%%%%%%%%%%%%%%%%%%%%%%%%%%%%%%%%%%%%%%%%%%%%

The $\Gamma$ estimates are \emph{priors}, not final answers. Following \citet{llmpriors2025scientific},
they should be updated with empirical data via Bayes' rule:

\begin{equation}
  p(\gamma_{dj} | D) \propto p(D | \gamma_{dj}) \cdot p_{LLM}(\gamma_{dj})
\end{equation}

\paragraph{When Priors Dominate.}
LLM-derived priors dominate when data are sparse, noisy, or unavailable---precisely the conditions
where structural models face the estimation gap.

\paragraph{When Data Dominate.}
Strong empirical evidence overwhelms priors. The $\Gamma$ estimates are defeasible:
they shape inference when evidence is weak but yield to strong evidence.

\paragraph{The Workflow.}
\begin{center}
\fbox{\parbox{0.85\textwidth}{
\centering
\textbf{From Implicit Knowledge to Updated Estimates}\\[0.5em]
$\underbrace{p_{LLM}(\gamma)}_{\text{LLM-MC Prior}}$ $\times$ $\underbrace{p(D|\gamma)}_{\text{Empirical Likelihood}}$ $\propto$ $\underbrace{p(\gamma|D)}_{\text{Posterior}}$\\[0.5em]
\textit{Start with what the literature knows. Update with what you observe.}
}}
\end{center}

%%%%%%%%%%%%%%%%%%%%%%%%%%%%%%%%%%%%%%%%%%%%%%%%%%%%%%%%%%%%%%%%%%%%%%%%%%%%%%%
\subsection{Summary: The Complete FEPSDE-$\Psi$ Framework}
%%%%%%%%%%%%%%%%%%%%%%%%%%%%%%%%%%%%%%%%%%%%%%%%%%%%%%%%%%%%%%%%%%%%%%%%%%%%%%%

\begin{table}[htbp]
\centering
\caption{The Complete Welfare Framework}
\label{tab:complete-framework}
\renewcommand{\arraystretch}{1.4}
\begin{tabular}{@{}llll@{}}
\toprule
\textbf{Element} & \textbf{Symbol} & \textbf{Source} & \textbf{Question Answered} \\
\midrule
Coherence & $K$ & Ch. 5--7 & Are we in equilibrium? \\
Quality & $Q$ & This chapter & Is the equilibrium good? \\
Context & $\Psi$ & App. V & What environment shapes welfare? \\
Interactions & $\Gamma$ & LLM-MC (App. AN) & How does context shape welfare? \\
Dimensions & FEPSDE & This chapter & Good in what respects? \\
Categories & INU/KNU/IDN & This chapter & Good for whom? \\
Time & $t = 0,1,2,3$ & This chapter & Good when? \\
Valence & G, P & Prospect Theory & Gains or avoided pains? \\
\midrule
\multicolumn{3}{l}{\textbf{Total Components:}} & $3 \times 6 \times 4 \times 2 = 144$ (+4 IDN = \textbf{148}) \\
\bottomrule
\end{tabular}
\end{table}

The FEPSDE-$\Psi$ framework provides the vocabulary to discuss welfare
beyond GDP, incorporating insights from psychology, health economics,
environmental economics, behavioral economics, and institutional economics
into a unified, context-sensitive structure.

\begin{center}
\fbox{\parbox{0.95\textwidth}{
\centering
\textbf{The Complete FEPSDE-$\Psi$ Model}\\[0.5em]
$Q(C^*, \Psi) = \displaystyle\sum_{i \in \{INU,KNU,IDN\}} \sum_{d \in FEPSDE} \sum_{t=0}^{3} \omega_i(\Psi) \cdot \delta_d(\Psi)^t \left[ G_{i,d,t}(\Psi) - \lambda_d(\Psi) \cdot P_{i,d,t}(\Psi) \right]$\\[0.5em]
where $U_d(\Psi) = \bar{U}_d + \sum_{j=1}^{8} \gamma_{dj} \cdot \Psi_j$\\[0.3em]
\textit{148 components $\times$ 8 context dimensions $\times$ 48 $\Gamma$-parameters}\\[0.3em]
\textit{LLM-MC validated: MAE = 0.034, sign agreement = 100\%}
}}
\end{center}

\vspace{1em}
\hrule
\vspace{0.5em}
\noindent\textit{Cross-references: Appendix~AN (LLM Monte Carlo methodology), 
Appendix~V ($\Psi$-Dimensions operationalization), 
Appendix~AS (Identity utility structure),
Appendix~C (Complete 144-component matrix with examples)}

%%%%%%%%%%%%%%%%%%%%%%%%%%%%%%%%%%%%%%%%%%%%%%%%%%%%%%%%%%%%%%%%%%%%%%%%%%%%%%%
